\documentclass[10.5pt,a4paper]{article}
\usepackage{fontspec}
\usepackage[french]{babel}
\usepackage[margin=1.5cm,top=1.2cm,bottom=1.2cm]{geometry}
\usepackage{fontawesome5}
\usepackage{titlesec}
\usepackage{enumitem}
\usepackage[hidelinks]{hyperref}
\usepackage{tabularx}
\usepackage{needspace}

\hypersetup{
  pdftitle={CV — Mohamed Amine Charrada},
  pdfauthor={Mohamed Amine Charrada},
  pdflang={fr}
}

\pagestyle{empty}
\setlength{\parindent}{0pt}
\frenchsetup{StandardLists=true}

\titleformat{\section}{\large\bfseries}{}{0pt}{}[\vspace{-6pt}\rule{\textwidth}{0.6pt}]
\titlespacing*{\section}{0pt}{8pt}{5pt}

\newcommand{\entry}[2]{\noindent #1 \hfill #2\par}
\newcommand{\project}[1]{\needspace{5\baselineskip}\noindent\textbf{#1}\par\vspace{2pt}}
\setlist[itemize]{leftmargin=1.2em,itemsep=0pt,topsep=3pt,label=\textbullet}
\newenvironment{tasks}{\begin{itemize}[leftmargin=2.6em,topsep=3pt]}{\end{itemize}\vspace{3pt}}

\begin{document}

\begin{center}
  {\LARGE\bfseries MOHAMED AMINE CHARRADA}\\[4pt]
  Étudiant en cycle ingénieur informatique\\[3pt]
  \faPhone*\ +216 52 904 114 \quad
  \faEnvelope\ \href{mailto:mohamedamine.charrada@esprit.tn}{mohamedamine.charrada@esprit.tn} \quad
  \faMapMarker*\ Tunis, Tunisie\\
  \faLinkedin\ \href{https://www.linkedin.com/in/mohamed-amine-charrada-17a494315/}{LinkedIn : Mohamed Amine Charrada} \quad
  \faGithub\ \href{https://github.com/charrada-amine}{GitHub : charrada-amine}
\end{center}

\section*{Profil}
Élève ingénieur en 5\textsuperscript{e} année à ESPRIT, orienté full stack, cloud et IA. Expérience de bout en bout, du backend
aux interfaces mobiles, web et desktop, jusqu’au déploiement cloud et CI/CD, avec un fort souci de la qualité logicielle.

\section*{Formation}
\entry{\textbf{ESPRIT} — Cycle ingénieur en informatique}{2022 – aujourd’hui}

\section*{Expériences professionnelles}
\entry{\textbf{Stagiaire en développement full stack — MEDCOM} (hybride)}{Juin 2026 – Oct. 2026}
\begin{tasks}
  \item Stage de 4 mois sur Zennyt : mini-jeux d’évaluation de bout en bout (Flutter et Spring Boot).
  \item Développement des API backend, des tests automatisés et de la console d’administration.
\end{tasks}
\entry{\textbf{Stagiaire en développement logiciel — ExpressDisplay} (sur site)}{Juin 2025 – Juil. 2025}
\begin{tasks}
  \item Développement de deux modèles de prédiction du temps en Python.
  \item Applications mobiles de reporting et de gestion RH (React Native).
\end{tasks}
\entry{\textbf{Stage d’intégration en réseaux et cybersécurité — OACA} (sur site)}{Juil. 2023 – Août 2023}
\begin{tasks}
  \item Étude de l’architecture réseau et des équipements d’infrastructure, et sensibilisation à la cybersécurité.
\end{tasks}

\section*{Projets}
\project{Zennyt — Plateforme de mise en relation entre candidats et recruteurs}
\begin{itemize}
  \item Application mobile, web et desktop : matching par swipe et Fit Score, mini-jeux d’évaluation anti-triche.
  \item Flutter, Spring Boot (Java 21), PostgreSQL, FastAPI, Docker ; hébergé sur Azure Container Apps.
\end{itemize}
\project{SamOps — Plateforme d’observabilité et d’optimisation des coûts cloud}
\begin{itemize}
  \item Monitoring, dashboards coûts/performance, CI/CD et agents d’analyse de la qualité du code.
  \item Microservices (Go, NestJS, Python) ; déployé sur AWS (\href{https://samops.app}{samops.app}).
\end{itemize}
\project{Stage — Plateforme de recherche de stage assistée par des agents IA (PWA)}
\begin{itemize}
  \item Agent de veille quotidien, matching par IA, coach d’entretien vocal temps réel, extension Chrome LinkedIn.
  \item Next.js 16, Convex, Gemini (Vercel AI SDK), Better Auth.
\end{itemize}
\project{DinoBot — Harness d’orchestration d’agents IA de développement (macOS)}
\begin{itemize}
  \item 10 agents spécialisés (développeur, testeur QA, DevOps, critique…) via Claude Code, Codex et API LLM.
  \item Runtime optimisé : sessions persistantes, cache de contexte, sandbox par agent ; Electron, React, Node.js.
\end{itemize}
\project{Styleto — Application mobile Smart Fashion (iOS / Android)}
\begin{itemize}
  \item Garde-robe numérique avec IA : recommandations, essayage virtuel, marketplace ; SwiftUI, Kotlin Compose, NestJS.
\end{itemize}

\section*{Compétences}
\begin{tabular}{@{\hspace{0.6em}}l}
  \textbf{Langages :} Java, Python, TypeScript, JavaScript, Go, Dart, Kotlin, Swift, C/C++, PHP, SQL\\
  \textbf{Frameworks :} Spring Boot, NestJS, Next.js, React, Flutter, React Native, FastAPI, Symfony\\
  \textbf{Bases de données :} PostgreSQL, MySQL, MongoDB, Firebase, Supabase\\
  \textbf{DevOps/Cloud :} Git, Docker, CI/CD, AWS, Azure, GCP, Linux
\end{tabular}

\section*{Certifications}
Hedera Certified Developer Associate \textbullet{} Microsoft Applied Skills : Azure Container Apps \textbullet{} AWS Cloud Foundations

\section*{Langues \& activités}
Arabe : maternelle \quad Français : B2 \quad Anglais : B2 \quad\textbar\quad IEEE (2023) \quad Croissant-Rouge tunisien (depuis 2021)

\end{document}
